\section{With the sources known}
\label{sec:truth}

\begin{figure*}[t]
  \centering
  \includegraphics[width=\textwidth]{fig1_axes.pdf}
  \caption{Held-out recovery along five EEG-relevant stress axes; harder to the right in every panel. Top row: permutation MCC against the signed sources (PCA, FastICA, iVAE, PCL), with the rotation reference (dashed), the score of a method that recovers the source subspace but not the sources. Bottom row: linear-invariant CCA-MCC against $|s|$ for the energy-like features of TCL, TCL after per-segment demeaning of the observations (dashed) and CEBRA-Time. Grey band: measured chance level (95th percentile of surrogate components). Error bars: bootstrap 95\% intervals over 20 seeds (PCA, FastICA, TCL), 5 seeds (iVAE, PCL, demeaned TCL, panel d, $M=32$) or 3 seeds (CEBRA). (a) $1/f$ background; (b) pointwise mixing nonlinearity; (c) first-order source coupling; (d) co-modulation fraction of the segment-wise modulation, from source-specific profiles to a single shared gain, with the range of nine BCI IV-2a subjects shaded; (e) number of sources against 22 channels, each $M$ with its own trained generator and learned mixing. Numbers in Table~\ref{tab:full}.}
  \label{fig:axes}
\end{figure*}

Figure~\ref{fig:axes} and Table~\ref{tab:full} give the result.
On clean linear mixing every method except PCA recovers the sources; PCA sits on the rotation reference everywhere, it finds the subspace and never separates it, which is why it serves as the subspace reference.
Pointwise nonlinearity and source coupling, the two axes closest to the classical theory, cost the nonlinear methods little.
The breaks are elsewhere, and each has a mechanism.

\paragraph{Drift feeds the segment classifier.}
As the $1/f$ background grows, TCL falls to the chance band while FastICA is barely affected (Figure~\ref{fig:axes}a).
The failure is not noise as such.
Power below the segment rate shifts the mean of each segment, the classifier learns segment identity from those shifts, and its features track drift rather than source variance; removing the per-segment mean of the observations restores TCL to its clean level down to 10~dB (dashed line, Appendix~\ref{app:robust}).
PCL fails from the other side: $1/f$ noise is temporally dependent, so it feeds a dependence-based objective, and PCL's own pair-discrimination accuracy rises along the axis while its recovery falls.
CEBRA, whose positive pairs are short time offsets, inherits TCL's profile without any labels.

\paragraph{A shared modulation identifies only total power.}
Along axis (d) every nonlinear method declines as the co-modulation fraction approaches one, TCL to the level the rank argument of Appendix~\ref{app:tcl} predicts: with one shared profile the modulation matrix has rank one, only $\|s\|^2$ is identified, and the individual sources collapse while the group energy is still decoded from the features (Appendix~\ref{app:robust}).
FastICA degrades too, since a common gain makes near-Gaussian sources spherically symmetric.
More data does not move these cells, nor the stationary one; the failure is in the assumption, not in the sample (Figure~\ref{fig:robust}b).
The ordering holds down to the modulation magnitudes of real EEG at second-scale segmentation (Figure~\ref{fig:real}a).

\paragraph{Sources against channels.}
As $M$ approaches $D$ the nonlinear methods fall towards the rotation reference, and past $D$ every method falls (Figure~\ref{fig:axes}e).
For TCL this is a finite-sample failure: with sixteen times more data its score at $M = D$ rises to that of the clean cell (Figure~\ref{fig:robust}b).
For the linear methods the limit is physical rather than statistical: a sphere-model leadfield in place of the learned mixing reproduces every ordering of Figure~\ref{fig:axes} and steepens this decline through its conditioning (Figure~\ref{fig:robust}d, Appendix~\ref{app:robust}).
